% GENERATED FILE. Edit canonical lessons, then run npm run generate:books.
% canonical-source-hash: fnv1a64:128397b1d9fad8b5
% canonical-lessons: AR-W04-arbaa-sutur, AR-R45-the-last-page

\chapter{Writing --- Four Lines, And The Joiner That Is Not A Word}
\label{ch:ar-arbaa-sutur}
\hlchaptermodality{45}

\begin{chapteropening}
\textbf{By the end of this chapter:} \emph{I can write four connected lines of Arabic to a named reader, attaching the and-joiner to the following word and giving a reason with the because-word.}

Complete AR-W04-arbaa-sutur: four lines to one reader, at least one joined with the attached and-letter and one reasoned, then check no space follows it.
\end{chapteropening}

\section[wa]{Four lines, and the joiner that is not a separate word}

\label{lesson:AR-W04-arbaa-sutur}



\begin{quote}
The last lesson asked for writing \textbf{finished on time}. This one asks for writing that \textbf{holds together}.

Pick one person who will read it.
\end{quote}

\subsection*{What to know first}
Connected writing is not \textbf{longer} writing. Four short lines can be connected; one long line can fail to be. What makes it connected is that a sentence stops standing on its own.

You hold three joiners. \textbf{\ar{و}} adds. \textbf{\ar{لكن}} sets against. \textbf{\ar{لأن}} gives the reason.

And the first of those does something the others do not. \textbf{\ar{لكن}} and \textbf{\ar{لأن}} stand as words, with space on both sides. \textbf{\ar{و}} does not. It is written \textbf{onto the front of the next word}, with no space at all.

So in Arabic the join is not a word the eye can skip over on its way to the content. It is a letter the following word now begins with — the two things being joined are physically fastened, and a reader meets the joint before they meet the second thing.

A hand that leaves a space after \textbf{\ar{و}} has not made a small typographic slip. It has written a letter that now belongs to nothing.

\subsection*{Writing — connected composition}
\hlblockfigure{\detokenize{figures/AR-W04-arbaa-sutur-filmstrip.pdf}}{How \ar{و} is written, stroke by stroke}

Write \textbf{four lines} to the person you picked, right to left.

At least one must use \textbf{\ar{و}} — attached, no space. At least one must use \textbf{\ar{لأن}}.

Use only words you already write. This lesson adds no vocabulary; the whole exercise is the joining.

\begin{tcolorbox}[breakable,colback=teal!4,colframe=teal!35!black,title={Before you move on}]
Cover your first line and read the second alone. Does it still make sense without the first? If it does, you wrote a \textbf{list} — which is where you were one lesson ago.

Then look at every \textbf{\ar{و}} you wrote. Is there a space after it? There should not be.
\end{tcolorbox}
\section[khayr · hamza · karrir min faḍlik]{The last page}

\label{lesson:AR-R45-the-last-page}



\begin{quote}
Write \textbf{\ar{خير}} by hand, right to left, and say the one sentence that gets a lost conversation moving. Both were taught once, at opposite ends of this book.
\end{quote}

\subsection*{Writing: the hook family, and the word it builds}
\textbf{\ar{ح} \ar{ح} \ar{خ} \ar{ج}} share one body: a rounded \textbf{hook} flowing into a \textbf{tail that drops below the line}. Only the dot tells them apart — \textbf{\ar{ح}} has none, \textbf{\ar{خ}} takes one above, \textbf{\ar{ج}} one below.

Write these out:

\begin{itemize}
\raggedright
  \item \textbf{\ar{خير}} — \emph{khāʾ} with its dot above, reaching left into \emph{yāʾ}, then \emph{rā}, and lift
  \item \textbf{\ar{بخير}} — glue \textbf{\ar{بـ}} on the front, one dot below, joining into the \emph{khāʾ}
\end{itemize}

Nothing joins after \textbf{\ar{ر}}, which is why the pen lifts. And \textbf{\ar{خير}} is the word that was hiding inside the morning greeting before anyone told you it was there; adding \textbf{\ar{بـ}}, \emph{in}, turns \emph{good} into the answer \textbf{in goodness}.

\subsection*{Script: the mark that is not a vowel}
Say English \emph{uh-oh} slowly and feel the catch between the halves. That catch is a \textbf{glottal stop}, and Arabic writes it as a real consonant: \textbf{\ar{ء}}, the \textbf{hamza}. It is not one of the short-vowel marks. Sitting on an \emph{alif} seat it becomes \textbf{\ar{أ}}, which is how the word for \emph{you} begins.

The short-vowel marks do something else entirely, and the clearest proof is the word for \emph{you}: \textbf{\ar{أنت}} is one consonant skeleton, and a single mark that ordinary print omits decides whether you have addressed \textbf{\ar{أنتَ}}, a man, or \textbf{\ar{أنتِ}}, a woman.

\subsection*{Your turn: the sentence to keep}
\begin{quote}
\textbf{\ar{كرّر} \ar{من} \ar{فضلك}} — \emph{karrir min faḍlik}
\end{quote}

Say these aloud:

\begin{itemize}
\raggedright
  \item you have not understood, then ask for the sentence back — \emph{lā afham. karrir min faḍlik}
  \item call somebody first, then ask — \emph{yā …, karrir min faḍlik}
\end{itemize}

Of everything in this book, that is the line to carry out of it. The marks Arabic leaves off the page are the ones a beginner most needs, so a beginner will miss words — and the fix is not a better memory. It is four syllables asking for another pass.

\begin{tcolorbox}[breakable,colback=teal!4,colframe=teal!35!black,title={Before you move on}]
Say whether hamza is a vowel mark and what it writes. (\textbf{No} — a glottal stop.) Say what tells \textbf{\ar{ح}}, \textbf{\ar{خ}} and \textbf{\ar{ج}} apart. (\textbf{The dot} — none, one above, one below.) Say what \textbf{\ar{بـ}} adds to \textbf{\ar{خير}}. (\emph{\textbf{In}} — in goodness.) And give the sentence that rescues a conversation. (\emph{\textbf{Karrir min faḍlik}}.)
\end{tcolorbox}
